\documentclass[11pt]{article}

\usepackage[a4paper,margin=1in]{geometry}
\usepackage{amsmath}
\usepackage{amssymb}
\usepackage{booktabs}
\usepackage{hyperref}
\usepackage{xcolor}
\usepackage{listings}

\lstset{
  basicstyle=\ttfamily\small,
  breaklines=true,
  frame=single,
  columns=fullflexible
}

\title{MCGDF\_object\_mlp:\\Robot + Object World Model with a Pure-MLP Object Side}
\author{}
\date{}

\begin{document}
\maketitle

\section{Overview}

The implementation under
\begin{center}
\texttt{scripts/world\_model/Physics/Robot\_and\_Object/MCGDF\_object\_mlp}
\end{center}
is a deliberate fork of the MCGDF model in the sibling folder
\texttt{MCGDF/} with one structural change on the object side:
\textbf{the Newton--Euler object integrator is removed entirely}.  In
its place, a single MLP head (\texttt{ObjectStateMLP}) emits the
13-dimensional next object substate directly from the current
$(s_t, \tau_t)$, with FiLM modulation by the variational latent
context $z$.

The robot side is unchanged from MCGDF: DeLaN-structured inertia /
Coriolis / gravity terms; an implicit-PD residual head $r_\theta$; the
joint damping / friction parameters; the analytic Franka FK +
geometric Jacobian; and the action--reaction torque
$\tau_{\mathrm{contact}} = -J_c^\top F_c$ that flows from a
still-learned contact wrench $F_c$.  All of those are preserved.

This document describes the architecture, the precise role the
contact wrench $F_c$ plays now that it no longer enters an object
integrator, the design space the new head explores, the impact on
training losses, and the asymmetric physical inductive bias compared
to its three sibling variants.

\section{State, Input, and Context}

State and input are identical to MCGDF:
\begin{equation}
  s_t =
  \begin{bmatrix}
    q_t \\ \dot q_t \\ p^o_t \\ r^o_t \\ v^o_t \\ \omega^o_t
  \end{bmatrix} \in \mathbb{R}^{18 + 13} = \mathbb{R}^{31},
  \qquad
  \tau_t \in \mathbb{R}^{9}.
\end{equation}
The per-episode physical context $\xi = [m_o,\,I_b\,\mathrm{(9)},\,
\mu_o\,(3)] \in \mathbb{R}^{13}$ is still loaded by the dataset and
still consumed by the \texttt{ContactWrenchHead} (so $F_c$ knows the
nominal physical context); it is \emph{not} consumed by the new
\texttt{ObjectStateMLP}.  The variational latent $z$ is produced by
the existing \texttt{ContextEncoder} from the $K$-step history window
exactly as in MCGDF.

\section{Forward Dynamics}

\subsection{Robot side (unchanged from MCGDF)}

\begin{equation}
  H(q,z)\,\ddot q
  \;=\;
  \tau + r_\theta(q,\dot q,\tau,z)
  - c(q,\dot q,z) - g(q,z)
  - d\dot q - f\tanh(\dot q/\varepsilon_f)
  - J_c(q)^\top F_c(s,\tau,\xi,z),
  \label{eq:om-robot}
\end{equation}
followed by semi-implicit Euler for the robot kinematics:
\begin{equation}
  \dot q_{t+1} = \dot q_t + \ddot q_t\,\Delta t,
  \qquad
  q_{t+1} = q_t + \dot q_{t+1}\,\Delta t.
\end{equation}
Eq.~\eqref{eq:om-robot} is identical to MCGDF.  The contact wrench
$F_c \in \mathbb{R}^6$ is still predicted by
\texttt{ContactWrenchHead}; it is still applied to the robot side via
Newton's third law.  What changes is that $F_c$ \emph{does not enter}
the object integrator anymore.  The MLP-based object head
(Section~\ref{sec:om-object}) is the sole determinant of
$(p^o_{t+1}, r^o_{t+1}, v^o_{t+1}, \omega^o_{t+1})$.

\subsection{Object side: a single MLP}
\label{sec:om-object}

The Newton--Euler equations
$m_o \dot v^o = m_o g + F_c^{\mathrm{lin}}$ and
$I_w \dot \omega^o + \omega^o\!\times\!I_w\omega^o = F_c^{\mathrm{ang}}$
are \emph{removed}.  In their place, a single MLP head
\texttt{ObjectStateMLP} emits the next 13-D object substate
directly:
\begin{equation}
  \hat o_{t+1}
  \;=\;
  \mathrm{ObjectStateMLP}(s_t,\,\tau_t,\,z)
  \;\in\; \mathbb{R}^{13}.
  \label{eq:om-objmlp}
\end{equation}
The 13 output dimensions are interpreted slot-by-slot as
\begin{equation}
  \hat o_{t+1}
  \;=\;
  \bigl(
    \hat p^o_{t+1}\,(3),\,
    \tilde r^o_{t+1}\,(4),\,
    \hat v^o_{t+1}\,(3),\,
    \hat \omega^o_{t+1}\,(3)
  \bigr),
\end{equation}
where the quaternion slot $\tilde r^o_{t+1}$ is normalized to unit
length before being written into the next state:
\begin{equation}
  r^o_{t+1} \;=\; \frac{\tilde r^o_{t+1}}{\|\tilde r^o_{t+1}\| + \varepsilon_q}.
\end{equation}
This single normalization step is the entire amount of analytic
geometry that survives on the object side; there is no
exp-map quaternion integrator, no semi-implicit Euler on
$(v^o, p^o)$, and no inertia or mass anywhere.

\subsection{ObjectStateMLP architecture}

\texttt{ObjectStateMLP} reuses the same FiLM pattern as
\texttt{ResidualTorqueHead} and \texttt{ContactWrenchHead}:
\begin{lstlisting}
in   = cat(state, torque)         # in R^{state_dim + torque_dim} = R^{40}
h    = SiLU(Linear_in(in))        # in R^{hidden_dim}        (= 256 default)
for i in 1..depth:                # depth FiLM blocks         (= 3 default)
    h = FiLM_i(h, z)              #   h <- SiLU(gamma_i(z) * Linear_i(h) + beta_i(z))
out  = Linear_out(h)              # in R^{13}
\end{lstlisting}
Each \texttt{FiLM\_i} is the existing \texttt{FiLMLayer}: an inner
linear projection whose output is element-wise affine-modulated by
\((\gamma_i(z), \beta_i(z))\) from two small linear maps of $z$, then
passed through SiLU.  When the context encoder is off, the head
gracefully falls back to a plain MLP without FiLM
(\texttt{latent\_dim is None}).

\paragraph{Two new CLI flags.}
\texttt{--object\_mlp\_hidden\_dim} (default $256$) and
\texttt{--object\_mlp\_depth} (default $3$) control the head's size.
They are saved into the checkpoint config and rehydrated by the
checkpoint loader in \texttt{plot\_multistep\_trajectory.py} and
\texttt{animate\_episode.py}.

\section{Why $F_c$ Is Still Predicted}

The contact wrench $F_c \in \mathbb{R}^6$ from
\texttt{ContactWrenchHead} no longer drives the object dynamics, but
it is \emph{not} dropped.  Three reasons:

\begin{enumerate}
  \item \textbf{Robot-side action--reaction (Newton's third law).}
    The robot ODE
    Eq.~\eqref{eq:om-robot} still relies on
    $\tau_{\mathrm{contact}} = -J_c^\top F_c$.  If the cube is being
    carried, there must be an upward reaction torque on the gripper;
    if the cube resists motion, there must be a backward reaction
    torque.  Predicting that reaction with a structured wrench (rather
    than absorbing it into $r_\theta$) keeps the robot-side dynamics
    interpretable and avoids the identifiability collapse documented
    in the MCGDF description's residual section.
  \item \textbf{Continuity with the rest of the codebase.}  Existing
    losses
    ($\mathcal{L}_{\mathrm{contact\,L2}} = \lambda_{\mathrm{contact}}\|F_c\|^2$,
    the auxiliary $F_c$ diagnostics in
    \texttt{aux["contact\_wrench"]} / \texttt{aux["tau\_contact"]}),
    and the existing animation / signal-panel overlays all keep
    working without modification.
  \item \textbf{Optional future re-coupling.}  Keeping $F_c$ alive
    leaves the door open to feeding $\|F_c^{\mathrm{lin}}\|$ as an
    extra input to the \texttt{ObjectStateMLP} (matching the
    coupling-gate trick in \texttt{No\_object\_Euler}), if you ever
    want to re-introduce a soft physical link.  Today $F_c$ flows
    only to the robot side.
\end{enumerate}

\section{Auxiliary Outputs and Loss Compatibility}
\label{sec:om-aux}

The step's \texttt{aux} dict still emits keys downstream consumers
expect, so the training script and evaluation tooling continue to
work \emph{without} edits to the loss code.  The semantics of two
families of keys changes, however, and is documented here.

\paragraph{Robot-side keys.}  Unchanged.  $H,\,c,\,g,\,
\ddot q,\,r_\theta,\,F_c,\,\tau_{\mathrm{contact}}$, and the
inverse-dynamics composite \texttt{full\_inverse\_dynamics\_tau} are
still produced exactly as in MCGDF.

\paragraph{Object-side keys.}  Re-interpreted as
\emph{finite-difference diagnostics}, not physical quantities:
\begin{align}
  \texttt{object\_lin\_acc}
    &= (\hat v^o_{t+1} - v^o_t) / \Delta t,
    \label{eq:om-fd-lin}\\
  \texttt{object\_ang\_acc}
    &= (\hat \omega^o_{t+1} - \omega^o_t) / \Delta t,
    \label{eq:om-fd-ang}\\
  \texttt{object\_gravity\_acc}
    &= (0,\,0,\,-9.81),\;\text{constant placeholder},\\
  \texttt{object\_inertial\_force},\;
  \texttt{object\_gravity\_force}
    &= \mathbf{0}\;\text{(no mass is available to scale these)}.
\end{align}
Eqs.~\eqref{eq:om-fd-lin}--\eqref{eq:om-fd-ang} are FD-derived
accelerations on top of the MLP's predicted velocities; they let
\texttt{supervised\_object\_dynamics\_loss} keep producing a signal
(it MSEs them against the recorded $\dot v^o_w / \dot \omega^o_w$),
which in turn constrains the MLP to produce kinematically consistent
velocity changes.  They are \emph{not} physical accelerations -- in
particular they do not satisfy $F = m a$, since $F_c$ never entered
this calculation.

Two new aux keys are emitted as a convenience:
\begin{itemize}
  \item \texttt{object\_state\_pred\_raw}: the raw 13-D MLP output
    before quaternion normalization.  Useful for monitoring whether
    the quaternion slot's pre-normalization norm collapses near 0
    (which would make the post-normalization quaternion noisy).
  \item \texttt{object\_quat\_next\_raw}: the unnormalized quaternion
    slot \(\tilde r^o_{t+1}\), for the same monitoring purpose.
\end{itemize}

\section{Trajectory Prediction (rollout)}

A predicted trajectory over horizon $H$ is produced exactly as in
MCGDF, by chaining $H$ calls to the per-step dynamics with the latent
$z$ frozen across the rollout:
\begin{equation}
  s_{t+h+1} \;=\; \mathrm{RobotObjectMCGDFStep}\bigl(s_{t+h},\,\tau_{t+h+1},\,\xi,\,z\bigr),
  \qquad h = 0, 1, \dots, H-1.
  \label{eq:om-rollout}
\end{equation}
The recursion is the standard Python \texttt{for} loop in
\texttt{MultiStepRobotObjectMCGDFWorldModel.forward}; no changes
relative to MCGDF were needed there.

\paragraph{Where the learned signal flows.}
At step $h$ of the rollout, the predicted object substate at step
$h+1$ is entirely determined by the 13-D output of the
\texttt{ObjectStateMLP} evaluated at $(s_{t+h}, \tau_{t+h+1}, z)$.
The robot substate at step $h+1$ is still co-driven by $F_c$ via the
DeLaN ODE on the robot side.  So in this variant the object trajectory
has a 13-D learned bottleneck (vs.\ MCGDF's 6-D contact-wrench
bottleneck, which then went through Newton--Euler), while the robot
trajectory is still mediated by the structured 6-D $F_c$.

\section{Sub-network Inventory}

\begin{table}[h]
\centering
\small
\begin{tabular}{p{0.28\linewidth} p{0.32\linewidth} p{0.18\linewidth} p{0.10\linewidth}}
\toprule
Sub-network (class)              & Input                                  & Output                       & Hidden \\
\midrule
\texttt{g\_net} (MLP)            & $q \in \mathbb{R}^9$                  & $g \in \mathbb{R}^9$         & 256    \\
\texttt{l\_diag\_net} (Deriv.\,MLP) & $q \in \mathbb{R}^9$              & $\ell_d \in \mathbb{R}^9$    & 256    \\
\texttt{l\_offdiag\_net} (Deriv.\,MLP) & $q \in \mathbb{R}^9$           & $\ell_o \in \mathbb{R}^{36}$ & 256    \\
\texttt{ResidualTorqueHead}      & $(q,\dot q,\tau) \in \mathbb{R}^{27}$ & $r_\theta \in \mathbb{R}^9$  & 128    \\
\texttt{ContactWrenchHead}       & $(s,\tau,\xi) \in \mathbb{R}^{53}$    & $F_c \in \mathbb{R}^6$       & 128    \\
\texttt{ObjectStateMLP}          & $(s,\tau) \in \mathbb{R}^{40}$        & $\hat o_{t+1} \in \mathbb{R}^{13}$ & 256 \\
\texttt{ContextEncoder} (VAE)    & flat history of $(s,\tau)$ over $K$ steps & $(\mu,\log\sigma^2) \in \mathbb{R}^{2\,\mathrm{ld}}$ & 256 \\
\texttt{JointDampingFrictionParams} (opt.) & ---                         & $9+9$ scalars                & ---    \\
\bottomrule
\end{tabular}
\caption{Trainable sub-networks of the MCGDF\_object\_mlp variant.
The new row is \texttt{ObjectStateMLP}; all other heads are
inherited from MCGDF.  ld $=$ \texttt{latent\_dim}; default $=8$.}
\label{tab:om-subnets}
\end{table}

The whole object side is therefore a single MLP plus a 4-D
quaternion normalization.  No analytic integration, no inertia
tensor, no mass, no gravity, no gyroscopic correction.

\section{What Is Gained and Lost}

\paragraph{Gained.}
\begin{itemize}
  \item \emph{Expressivity.}  An unconstrained 13-D regressor can
    represent any one-step object motion that lies inside the
    network's capacity, including non-rigid-body effects (slip in the
    pinch, micro-vibration, finger compliance, contact transitions)
    that the original Newton--Euler integrator literally cannot
    express because they violate $F = m a$ at fixed $(m_o, I_b)$.
  \item \emph{No oracle physics at inference.}  The Newton--Euler
    integrator in MCGDF consumed $m_o$ and $I_b$ from the dataset's
    $\xi$.  This variant's \texttt{ObjectStateMLP} does not -- so a
    deployment-time evaluation where $\xi$ is unknown does not need a
    $\hat\xi(z)$ regression head for the object side (the
    contact-head consumer of $\xi$ is the only remaining channel
    where $\xi$ enters the forward pass).
  \item \emph{Architectural simplicity.}  The whole object side
    collapses to one MLP forward call.  No quaternion exp map, no
    Cholesky-style inertia, no inertia rotation.
\end{itemize}

\paragraph{Lost.}
\begin{itemize}
  \item \emph{All physical inductive bias on the object side.}  The
    model now has to learn $F = m a$, the gyroscopic correction, the
    persistence of gravity, the kinematic identity
    $\dot p = v$, and the rotation update law -- all from data, with
    no structural prior.  Out-of-distribution behaviour is therefore
    much harder to bound than in MCGDF.
  \item \emph{Quaternion-on-the-sphere structure.}  Pure-MLP
    regression to a 4-D vector + normalization does \emph{not}
    preserve the geodesic geometry of \(SO(3)\): the loss can have
    bad local minima near antipodal quaternions, the network can
    pick the wrong hemisphere, and small numerical errors can
    accumulate into a drift on the sphere across long rollouts.  The
    \texttt{weighted\_rollout\_mse} already uses a sign-invariant
    quaternion MSE which helps a little.
  \item \emph{Stiff coupling between $F_c$ and the cube is gone.}
    MCGDF's $F_c$ \emph{had to} produce the observed cube
    acceleration via $F = m a$; here, $F_c$ is decoupled from the
    cube on the object side, so there is no longer a structural
    constraint forcing $F_c$ to be physically meaningful.  Adding
    L2 regularization on $\|F_c\|$ remains essential.
  \item \emph{No analytic energy bound.}  MCGDF's Newton--Euler
    integrator preserves total kinetic energy up to the work done by
    $F_c$.  The pure-MLP integrator can produce trajectories with
    arbitrary energy creation / dissipation per step -- there is no
    structural guarantee.
\end{itemize}

\section{Comparison with the Sibling Models}

\begin{table}[h]
\centering
\small
\begin{tabular}{p{0.20\linewidth} p{0.34\linewidth} p{0.20\linewidth} p{0.16\linewidth}}
\toprule
Model & Object-side equation & Learned signal driving the object & Physical prior \\
\midrule
\textbf{MCGDF} & Newton--Euler:
\(m_o\dot v^o = m_o g + F_c^{\mathrm{lin}}\),
\(I_w \dot\omega^o + \omega^o\!\times\! I_w\omega^o = F_c^{\mathrm{ang}}\) +
semi-implicit Euler + exp-map quat
& Contact wrench $F_c \in \mathbb{R}^6$
& Strong: $F=ma$, gyroscopic Euler, world-frame inertia, quat-exp \\
\textbf{No-Object-Euler} & Kinematic blend:
\(v^o_{t+1} = \alpha_t v_g + (1-\alpha_t)\rho_v v^o_t\)
(symmetric for $\omega^o$) + exp-map quat
& Coupling factor $\alpha_t \in [0,1]$ (gated by $\|F_c^{\mathrm{lin}}\|$ and $z$)
& Medium: carry vs.\ release dichotomy + quat-exp \\
\textbf{Context DeLaN} & Acceleration head:
\((\dot v^o, \dot\omega^o) = \mathrm{MLP}(s,\tau,z)\) + semi-implicit Euler + exp-map quat
& The 6-D acceleration vector
& Low: integrators preserve kinematic identities, but no $F=ma$ \\
\textbf{MCGDF\_object\_mlp} (this) & Next-state head:
\(\hat o_{t+1} = \mathrm{MLP}(s,\tau,z)\), quaternion slot normalized
& The 13-D next-state vector itself
& \emph{None.}  No analytic integrator, no $F=ma$, no quat-exp. \\
\bottomrule
\end{tabular}
\caption{Object-side comparison across the four sibling models in
\texttt{Robot\_and\_Object/}.  Reading left-to-right inside the
``Physical prior'' column lays out the four variants on a
strong-to-none spectrum.}
\label{tab:om-vs-others}
\end{table}

This variant sits at the most flexible end of the spectrum: it makes
the weakest physical assumption on the object side of any of the four
sibling implementations.  That makes it the natural baseline for
``how far can a pure MLP go on this task'' and the natural
upper-bound capacity comparison against the structured
\texttt{MCGDF} and \texttt{No\_object\_Euler} variants.

\section{Practical Notes}

\paragraph{Training.}  No changes to losses are required.
\texttt{wm\_train.py} produces exactly the same loss surface as
MCGDF, including the supervised
$\dot v^o / \dot \omega^o$ MSE -- which now constrains the
finite-differenced velocities of the MLP rather than the analytic
Newton--Euler accelerations.  In our experience the rollout MSE on
position is the dominant signal for this variant; the supervised
acceleration losses can be down-weighted (e.g.\
\texttt{--lambda\_object\_lin\_acc 0} and
\texttt{--lambda\_object\_ang\_acc 0}) without major regression.

\paragraph{Evaluation.}  The existing
\texttt{plot\_multistep\_trajectory.py},
\texttt{animate\_episode.py}, \texttt{plot\_episode\_dynamics.py},
and \texttt{plot\_pred\_error\_heavy\_light.py} all keep working.
The animation script's signal panel still draws
\texttt{contact\_wrench\_lin} / \texttt{contact\_wrench\_ang} from
\texttt{aux} -- those panels now describe the robot-side $F_c$ only;
they no longer tell you anything about the object's acceleration.

\paragraph{Where to look in the code.}
\begin{itemize}
  \item \texttt{models.py}: new class \texttt{ObjectStateMLP};
    modified \texttt{RobotObjectMCGDFStep.\_\_init\_\_} (adds the
    head), \texttt{RobotObjectMCGDFStep.forward} (replaces the
    Newton--Euler block with the MLP call + quaternion normalization
    + FD-diagnostic aux fields), and
    \texttt{MultiStepRobotObjectMCGDFWorldModel.\_\_init\_\_} (passes
    the two new size flags through).
  \item \texttt{wm\_train.py}: new \texttt{TrainConfig} fields
    \texttt{object\_mlp\_hidden\_dim} and \texttt{object\_mlp\_depth},
    new CLI flags of the same names, plumbed into the model
    constructor.
  \item \texttt{plot\_multistep\_trajectory.py} and
    \texttt{animate\_episode.py}: their \texttt{load\_checkpoint\_model}
    helpers read the two new fields from \texttt{cfg} with sensible
    defaults so older MCGDF checkpoints can still be loaded
    side-by-side.
\end{itemize}

\end{document}
